% !TeX root=../main.tex
\chapter{تحلیل نظری و پیچیدگی محاسباتی}
\label{chap:complexity}

\section{مقدمه}
این فصل پیچیدگی فضای جست‌وجو، اثر قیود و نویز بر هزینهٔ محاسباتی، و کران‌های عملی روش پیشنهادی را به‌صورت مستقل از ارزیابی شبیه‌سازی تحلیل می‌کند.

\section{پیچیدگی فضای جست‌وجوی کامل}
ابعاد فضای کامل شامل ترکیب فضای پلی‌نوم، حافظه و جایگزشت‌ها توصیف می‌شود.

\section{پیچیدگی حالت brute-force}
جست‌وجوی کامل درهم‌ریز از مرتبهٔ فاکتوریل است و عملاً غیرقابل اجراست؛ این بخش کران پایین محاسباتی را روشن می‌کند.

\section{پیچیدگی جست‌وجوی محدودشده}
با اعمال قیود جبری و آماری، فضای مؤثر به‌صورت قابل توجهی کاهش می‌یابد؛ مدل پیچیدگی الگوریتم مقید در اینجا ارائه می‌شود.

\section{تحلیل Branching Factor}
ضریب انشعاب درخت جست‌وجو در اعماق مختلف و وابستگی آن به نوع قیود بررسی می‌شود.

\section{اثر قیود بر کاهش فضای جست‌وجو}
هر قید (پاریتی، تحقق‌پذیری پلی‌نوم، سازگاری درهم‌ریز) سهم جداگانه‌ای در هرس دارد که باید کمی‌سازی شود.

\section{رابطه پیچیدگی با نویز}
با افزایش نویز، تعداد فرضیه‌هایی که از هرس عبور می‌کنند افزایش می‌یابد و رفتار جست‌وجو به سمت رشد نمایی میل می‌کند؛ این پدیده از یافته‌های مهم تحلیل است.

\section{اثر حافظه کدگذار}
افزایش \(m\) هم فضای پلی‌نوم و هم طول خاتمه را بزرگ‌تر می‌کند و هزینه را بالا می‌برد.

\section{اثر طول interleaver}
طول بزرگ‌تر جایگزشت، فضای نامزدها و زمان بسط گره‌ها را افزایش می‌دهد.

\section{تحلیل زمان اجرا}
مدل زمان اجرا برحسب تعداد گره‌های بسط‌یافته، آزمون‌های آماری و عملیات هرس ارائه می‌شود.

\section{کران‌های عملی روش}
شرایطی که در آن روش از نظر زمان و حافظه عملی باقی می‌ماند، همراه با رژیم‌های ناکارآمد، جمع‌بندی می‌گردد.
